%%% ВЫБОР УСТРОЙСТВА %%%

\IfStrEqCase{\devicetype}{%
{ЮНИТ-М500-ДЗТ}{%
% код для этого типа
\def\Type{Т} 
\def\Header{Технологические защиты трансформатора} 
\def\fileExt{t}
\def\AbzazOne{
\item
Технологические защиты от повышения давления, температуры масла и обмоток предназначены для предотвращения повреждения трансформаторного оборудования при нарушениях режима работы. Для трансформатора    предусматриваются функциональные блоки приема и обработки сигналов от датчиков температуры масла и обмоток. 
\item
Алгоритм работы функций в составе ТЗ Т приведен на рисунке \ref{pic:TZ}. Уставки ТЗ Т приведены в таблице \ref{app_set:tz}.  
}
}%
{ЮНИТ-М500-РЗТ}{%
% код для этого типа
\def\Type{Т} 
\def\Header{Технологические защиты трансформатора}
\def\fileExt{t}
\def\AbzazOne{
\item
Технологические защиты от повышения давления, температуры масла и обмоток предназначены для предотвращения повреждения трансформаторного оборудования при нарушениях режима работы. Для трансформатора    предусматриваются функциональные блоки приема и обработки сигналов от датчиков температуры масла и обмоток. 
\item
Алгоритм работы функций в составе ТЗ Т приведен на рисунке \ref{pic:TZ}. Уставки ТЗ Т приведены в таблице \ref{app_set:tz}. 
}
}%
{ЮНИТ-М500-ДЗАТ}{%
% код для этого типа
\def\Type{АТ} 
\def\Header{Технологические защиты автотрансформатора}
\def\fileExt{at}
\def\AbzazOne{
\item
Технологические защиты от повышения давления, температуры масла и обмоток предназначены для предотвращения повреждения трансформаторного оборудования при нарушениях режима работы. Для автотрансформатора и ЛРТ предусматриваются функциональные блоки приема и обработки сигналов от датчиков температуры масла и обмоток; для ЛРТ также предусматривается прием сигналов от реле давления.
\item
Алгоритм работы функций в составе ТЗ АТ приведен на рисунке \ref{pic:TZ}, алгоритм работы ТЗ ЛРТ выполнен аналогично. Уставки ТЗ АТ приведены в таблице \ref{app_set:tzat}, уставки ТЗ ЛРТ приведены в таблице \ref{app_set:tzlrt}.
}
}%
}
[
\def\TypeT{??}
\def\HeaderT{??}
]


\color{unidarkgreen}{\subsection{\Header}}
\color{black}

\begin{enumerate}

\AbzazOne

\item
Предусмотрена возможность приема сигналов срабатывания сигнального и аварийного контактов датчиков температуры масла и обмотки автотрансформатора (ДТм \Type~и ДТо \Type), а также сигналов низкой изоляции цепей аварийных контактов, по двум независимым цепям от источников ПДС1 и ПДС2. Для приема сигналов от ПДС1 предусмотрены входные сигналы \textit{«Ср. сигн контакта ДТм (ДТо) \Type~1 цепь»}, \textit{«Ср. авар контакта ДТм (ДТо) \Type~1 цепь»} и \textit{«Низкая изол. ДТм (ДТо) \Type~авар 1 цепь»}; для ПДС2 — соответственно \textit{«... 2 цепь»}. Если сигналы заводятся непосредственно со шкафа ШОМ, конфигурируются только сигналы по первой цепи, а действие по второй цепи сигнального и аварийного контактов выводится уставками \textbf{«ДТм (ДТо) \Type~сигн 2 цепь»} и \textbf{«ДТм (ДТо)\Type~авар 2 цепь»}

\item
Ввод в работу ДТм \Type~(ДТо \Type) осуществляется с помощью уставок \textbf{«ДТм \Type~(ДТо \Type) сигн 1 цепь»}, \textbf{«ДТм \Type~(ДТо \Type) сигн 2 цепь»}, \textbf{«ДТм \Type~(ДТо \Type) авар 1 цепь»} и \textbf{«ДТм \Type~(ДТо \Type) авар 2 цепь»}. Предусмотрена возможность действия срабатывания аварийного контакта ДТм \Type~(ДТо \Type) на отключение трансформатора с помощью уставки \textbf{«Откл от ДТм \Type~(ДТо\Type) авар»}.

\item
Оперативный ввод/вывод ТЗ \Type~производится с помощью БУ (блока управления функциями) \textit{«Вывод ТЗ \Type»}.

\item
Непрерывный контроль изоляции цепи аварийного контакта ДТм (ДТо) осуществляется с помощью реле контроля изоляции (РКИ). В случае снижения сопротивления изоляции РКИ подает сигнал \textit{«Низкая изол. ДТм \Type~(ДТо \Type) авар 1 цепь»} (\textit{«Низкая изол. ДТм \Type~(ДТо \Type) авар 2 цепь»}) для блокировки срабатывания ДТм \Type~(ДТо \Type) от аварийного контакта. Время срабатывания блокировки от сигналов РКИ задается с помощью уставки \textbf{«Тбл ТЗ \Type»}. Ввод блокировки ДТм \Type~(ДТо \Type) от РКИ осуществляется с помощью уставки \textbf{«Блок ТЗ \Type»}. Для сброса блокировки от РКИ, используется команда управления \textit{«Возврат ТЗ \Type»}.

\vspace{2mm}
\begin{figure}[H]
\centering
\includegraphics[width=1\textwidth,height=1\textheight,keepaspectratio]{\fbpath/06.02 ДЗАТ/041.1061 ТЗ АТ/img/TZ_\fileExt.pdf}
\captionof{figure}{Функциональная схема логики ТЗ \Type}\label{pic:TZ}
\end{figure}


\end{enumerate}